% Part: turing-machines
% Chapter: machines-computations
% Section: turing-machines

\documentclass[../../../include/open-logic-section]{subfiles}

\begin{document}

\olfileid[tr]{tur}{mac}{tur}
\olsection{Turing Makineleri}

\begin{explain}
Bir Turing makinesini oluşturan şeyin biçimsel tanımı soyut görünür, fakat
aslında basittir: Turing makinesinin işleyişini belirtmek için gereken bütün
bilgiyi tek bir matematiksel yapı içinde toplar. Bunlar (1)~makinenin
bulunabileceği durumlar, (2)~şeritte kullanılmasına izin verilen simgeler,
(3)~makinenin başlaması gereken durum ve (4)~makinenin yönerge kümesidir.
\end{explain}

\begin{defn}[Turing makinesi]
Bir \emph{Turing makinesi} $M$, aşağıdakilerden oluşan
$\langle Q,\Sigma,q_0,\delta\rangle$ dörtlüsüdür:
\begin{enumerate}
\item Sonlu bir \emph{durumlar} kümesi $Q$.
\item $\TMendtape$ ile $\TMblank$ simgelerini içeren sonlu bir
  \emph{alfabe} $\Sigma$.
\item Bir \emph{başlangıç durumu} $q_0\in Q$.
\item Sonlu bir \emph{yönerge kümesi}
  $\delta\colon Q\times\Sigma\pto
  Q\times\Sigma\times\{\TMleft,\TMright,\TMstay\}$.
\end{enumerate}
Kısmi $\delta$ fonksiyonuna $M$'nin \emph{geçiş fonksiyonu} da denir.
\end{defn}

\begin{explain}
Şeridin yalnızca bir yönde sonsuz olduğunu varsayıyoruz. Bu nedenle özel bir
$\TMendtape$ simgesini şeridin sol ucunun işareti olarak belirlemek
yararlıdır. Bu, Turing makinesi programlarının şeridin dışına çıkma
``tehlikesi'' içinde olduklarını anlamalarını kolaylaştırır. Bu simgenin hiç
üzerine yazılmadığını, yani $\delta(q,\TMendtape)$ tanımlıysa
$\delta(q,\TMendtape)=\tuple{q',\TMendtape,x}$ olduğunu varsayabilirdik.
Bazı ders kitapları bunu yapar; biz yapmıyoruz. Turing makinenizi kurarken
$\TMendtape$ simgesinin üzerine hiçbir zaman yazmamasına dikkat edebilirsiniz.
Üstelik üzerine yazmaya izin vermenin uygun bir esneklik sağladığı durumlar
vardır.
\end{explain}

\begin{ex}
\emph{Çift Sayı Makinesi:} Çift sayı makinesi biçimsel olarak
$\tuple{Q,\Sigma,q_0,\delta}$ dörtlüsüdür; burada
\begin{align*}
Q & = \{q_0,q_1\}\\
\Sigma & = \{\TMendtape,\TMblank,\TMstroke\},\\
\delta(q_0,\TMstroke) & = \tuple{q_1,\TMstroke,\TMright},\\
\delta(q_1,\TMstroke) & = \tuple{q_0,\TMstroke,\TMright},\\
\delta(q_1,\TMblank) & = \tuple{q_1,\TMblank,\TMright}.
\end{align*}
\end{ex}

\end{document}
